\chapter{开始使用 KSSOLV Toolbox}\label{ch:introduction}
\chapterlead{KSSOLV Toolbox 把原子结构、计算步骤和结果组织在同一个工作区中。可以从一个分子或晶体开始，检查它的几何形状，选择计算方法，再沿着工作流得到总能量、电子能级或其他结果。本章说明本书的阅读方法，以及完成一次计算时需要建立的几个基本认识。}

\section{从结构走向结果}
对一个材料体系提出问题，往往先于打开软件。例如，改变晶格常数会怎样影响总能量？一个缺陷会在晶体中引入怎样的局部环境？分子的初始键长经过优化后会如何变化？KSSOLV Toolbox 为这些研究提供结构准备和电子结构计算工具。它使用平面波基组和赝势描述电子体系，由 KSSOLV 完成 Kohn--Sham 方程的数值求解。

开始练习时，不必先掌握全部理论。应先认清三件事：计算对象是什么，软件将执行哪些步骤，最终的量回答了什么问题。以氢分子为例，结构给出两个氢原子的位置；自洽场计算在这些位置不变时求电子基态；结构优化进一步调整原子坐标，寻找能量较低的几何形状。这些步骤之间的区别，比一次输入多少参数更值得先弄清。

软件中的典型过程如下。
\begin{steps}
\item \term{准备结构。}导入文件，或使用建模工具创建结构。检查元素、原子数、晶格及周期性。
\item \term{建立工作流。}从模板开始，选择结构输入、计算设置、求解和绘图等节点。
\item \term{设置计算条件。}确认赝势、截断能、$k$ 点、收敛阈值和需要的结果。
\item \term{运行并检查。}查看输出消息、迭代过程和结束状态，判断求解是否完成。
\item \term{分析与保存。}阅读数值与图形，记录输入条件，将工作流和结果保存在项目中。
\end{steps}

这个顺序也适用于稍复杂的任务。表面吸附需要先建立有真空层的表面模型；能带计算需要先获得自洽电子密度；远程运行则在提交前选好计算位置。后续章节会把这些条件放回具体实例中说明。

\section{选择使用方式}
\subsection{独立应用程序}
独立版安装在本机，由 MATLAB Runtime 提供运行环境。它不要求安装完整 MATLAB。项目和结构文件通常保存在本机磁盘，图形窗口与文件选择窗口也使用桌面环境。对希望把项目与本机资料放在一起管理的读者，这是直接的入口。

\subsection{KSSOLV Toolbox Online}
Online 通过浏览器打开工作空间。浏览器负责显示与交互，计算在服务器上的会话中执行。本机可以上传结构文件，也可以下载项目与导出结果。云端工作目录与本机下载目录是两个位置：保存到云端后，要带走本机副本，还需要执行下载。

两种入口共享结构、工作流和结果的主要概念。本书在通用操作中使用一致的名称，并在保存、文件传输、会话退出、计算资源等有差别的位置分别说明。第一次使用 Online，可先阅读第\ref{ch:online}章，再返回工作区和计算教程。

\subsection{MATLAB 工具箱}
已安装 MATLAB 的读者还可使用工具箱版。当前版本要求 MATLAB R2025a 或更新版本。这一方式便于将图形操作与 MATLAB 脚本配合使用。本书的主要教程不要求读者编写 MATLAB 代码；附录提供小型复算脚本，供希望进一步理解输入和结果的读者使用。

\section{项目怎样组织工作}
\term{项目}是一次工作的容器，扩展名为 \file{.ks}。一个项目可以包含多个结构、多个工作流及其结果。可将它理解为一本实验记录：既收录研究对象，也保留计算的组织方式。

\term{结构}说明原子是什么、位于何处，以及周期结构的晶胞如何定义。三维视图便于检查几何，但屏幕中的球体大小、背景颜色和显示重复范围不是电子结构计算参数。

\term{工作流}说明步骤如何相接。节点代表某项任务，连线表示数据的传递。把结构与 SCF 计算相连，可以让求解器读取它；把计算结果接到收敛图节点，可以查看迭代过程。工作流本身也可保存成 \file{.wf} 文件，便于在其他项目中重用。

\term{数据与图形}保存在项目浏览器的结果分支中。数据可以用于后续计算，图形便于观察变化。仅保存一张图片，通常不足以重新分析结果；应同时保留项目与必要的原始数据。

\begin{notebox}{先区分“看见的结构”和“参与计算的结构”}
将晶体的显示重复设为 $2\times2\times2$，可以帮助观察相邻晶胞，但不会自动创建八倍原子数的计算模型。若要研究真实超胞中的空位或掺杂，应使用\ui{构建超胞}。这一区别将在建模教程中实际演示。
\end{notebox}

\section{本书的阅读路线}
\subsection{第一次完成计算}
先阅读安装、Online 和工作区章节，接着完成第\ref{ch:first}章的氢分子计算。练习涵盖结构导入、打开模板、设置基本参数、运行求解与保存结果。随后阅读计算基础，可以把已经见过的参数与物理含义联系起来。

\subsection{准备材料模型}
如果已经熟悉计算，但需要构建自己的材料，可从第\ref{ch:modeling}章进入。硅晶体、苯环和乙醇分别展示周期结构与分子的日常操作；第\ref{ch:advanced}章把这些方法延伸到表面、异质结、扭转双层、纳米结构和软物质，第\ref{ch:modeling-tools}章则说明具体工具的参数与使用方法。每次改变结构后，先核对原子数、晶胞和局部距离，再把模型交给计算工作流。

\subsection{整理与展示结果}
第\ref{ch:volume}章介绍体数据，第\ref{ch:results}章介绍数值和图形结果。阅读时应把结果与输入条件放在一起：总能量属于哪一个结构，能带采用哪条路径，等值面来自哪个数据集。这样得到的图才具有清楚的含义，也便于在课程报告中解释。

\section{教程的约定}
\ui{主页 → 结构 → 从文件导入结构}表示依次选择页面、按钮和菜单项。若菜单仍显示英文，本书保留英文名称。\file{silicon.cif}等名称表示文件；键盘上的按键写成 \ui{Enter}、\ui{Esc} 或 \ui{Ctrl/Cmd+Z}。后者表示 Windows/Linux 常用 Ctrl，macOS 常用 Command。

本书随附的 \file{examples} 目录包含教程文件。练习前将它复制到自己的课程目录，例如 \file{KSSOLV-Study}。建议为每一例建立独立子目录，使输入结构、项目与导出图形保存在同处。文件名可以体现对象和用途，例如 \file{H2-SCF.ks}、\file{Si-vacancy.ks}；这样在数周后重新打开时，也容易认出内容。

教程中的小型参数用于演示操作与解释结果。正式研究还应围绕目标物理量做收敛检查。本书会在相关章节说明如何改变一个参数、比较结果，并判断继续加密是否有意义，而不要求初学者在第一次练习时完成全部研究设置。

\section{从一个清楚的问题开始}
软件可以很快生成复杂的结构和工作流，而可靠的练习通常从简单对象开始。先完成一个已知的小体系，再更换材料；先确认自洽计算，再增加能带或态密度步骤；先保存一个清楚的项目，再尝试批量操作。每一步都应能用一句话说明目的。

阅读本书时，可以把关键操作的检查结果写进自己的实验记录。记录不必很长：原子数、晶胞、主要参数、是否收敛，以及一两条观察就足够。随着练习深入，这些简短记录会帮助把不同阶段的模型与结果联系起来。


\section{按问题寻找功能}

初次学习不必顺着全部菜单逐项尝试。先说清希望得到什么，再找到产生该结果的操作：有了结构文件，先学导入和检查；没有结构文件，可以使用数据库检索或建模；要计算基态，建立SCF；要看电子在倒空间怎样变化，选择路径能带或能带曲面；要研究激发，再准备基态空轨道与TDDFT输入。

第\ref{ch:workspace}章介绍数据库检索、快捷工具条和命令窗口，第\ref{ch:workflows}章列出14项工作流任务，并分别说明它们需要什么输入。第\ref{ch:results}章区分数据集、图形和工作空间变量，第\ref{ch:remote}章介绍对话提问、应用设置与完整MATLAB的远程作业。附录的任务速查可从目标结果反查正文；遇到陌生任务时，先看需要的输入，再回到对应练习。
